% = PARAGRAPH 1 - What deepfakes are and how they work
% = (sets up technical context before discussing consequences)

Gen AI is the primary tool used in deepfake technologies, to manipulate images, videos or audio to convincingly replace one person's likeness with another's~\citep{muam2026role}. This is, besides for entertainment purposes, often done for malicious purposes. Early deepfakes were easy to identify because they were not very realistic. However, advances in Gen AI have changed that. Now, deepfakes are harder to detect and that makes them a powerful tool for disinformation campaigns, for example in conflict zones~\citep{laws14060083,muam2026role}.

% = PARAGRAPH 2 - Where and how AI is being applied in conflict zones
% = satisfies: "how are where AI is applied"

Disinformation using Gen AI is not a theoretical concern. It has been actively deployed in real-world conflict zones, like Ukraine. The DFRLab report~\citep{osadchuk2024ai} documents multiple instances where deepfakes were used to spread false information. Examples include fabricated videos of political leaders making inflammatory statements or staged military actions. In the case of Ukraine, deepfakes have been used to create confusion, with one video\footnote{https://www.youtube.com/watch?v=X17yrEV5sl4} showing a fabricated speech by Ukrainian president, Volodymyr Zelenskyy, calling his own troops to surrender themselves.

% = PARAGRAPH 3 - Who the stakeholders are and how consequences are created 
% = satisfies: "why and how consequences are created / who are the stakeholders"

The consequences of deepfakes in conflict zones do not just affect the government and soldiers. Civilian populations can be misled by fabricated videos, which often leads to panic or misguided actions. Opposing governments may use deepfakes as propaganda tools to justify aggressive policies. International bodies like the UN may struggle to respond effectively when visual evidence is no longer trustworthy. Journalists face challenges in verifying sources and maintaining credibility. The general public may become increasingly skeptical of all media because trust in information sources has eroded~\citep{muam2026role,laws14060083}.

% = PARAGRAPH 4 - The ethical challenge and whether it is a new or old problem
% = satisfies: "why it is an ethical challenge / new or old problem"

Wartime propaganda and deception are not new phenomena. They have been part of human conflict for centuries. However, AI changes the scale, accessibility, and believability of disinformation. Deepfakes were easily discernible in the past, but now increasingly people are struggling to distinguish real from fake. This creates a never-before-seen ethical challenge because it undermines trust in visual evidence. As \citet{kenderdine2025deep} argue, deepfakes erode our cultural and epistemic trust in images and videos, which have traditionally been powerful tools for conveying truth.

% = PARAGRAPH 5 - The role of societal factors beyond AI itself
% = satisfies: "reflect if other societal factors are involved"

Pre-existing societal conditions such as political polarization, weakened media literacy, and fractured international trust amplify the impact of deepfakes. In highly polarised societies, people may be more likely to accept deepfakes that align with their beliefs and reject those that do not. Countries actively in conflict are inherently polarised, making them susceptible to deepfakes that reinforce existing narratives. Weakened media literacy means that many individuals lack the skills to critically evaluate sources. This makes them more susceptible to manipulation. Additionally, fractured international trust can hinder coordinated responses to disinformation campaigns.